\section{Final remarks}
\label{sec:conclusion}

We have presented \ours, a framework that transforms high-performing neural CATE estimators into knowledge-augmented networks with closed-form, auditable effect functions. On standard benchmarks, at least one \our variant matches or surpasses its neural counterpart on PEHE/ATE, while shallow heads (KAAMs) often yield the most favorable accuracy-interpretability trade-off. The pruning and auto-symbolification stages expose explicit formulas and partial dependence plots (PDP; \citet{friedman2001greedy}). These properties make \ours suitable for healthcare, policy evaluation, and economics, where interpretability is central to adoption and auditing \citep{AmannBMC2020, TonekaboniMLHC2019, Rudin2018StopEB, DoshiVelez2017TowardsAR}. Closed-form expressions can be scrutinized, simplified, or aligned with domain-grounded terms to support external validity. We release code to ensure reproducibility and facilitate adoption.

Nevertheless, \ours also present open challenges. Visualization tools are currently most effective for shallow KAAMs, while scaling them to deeper or more complex networks remains an area for development. Training and inference are generally more demanding than for standard neural models, and performance can be sensitive to regularization choices or pruning/symbolification budgets, which may introduce approximation errors or variability. As with all observational CATE estimation, our conclusions depend on standard identification assumptions (e.g., unconfoundedness, overlap), and our evaluation relies mainly on semi-synthetic datasets, so further assessment on field data is needed. Finally, uncertainty quantification for symbolic outputs and broader support for continuous or multi-valued treatments are still limited.

Building on these observations, future work should expand theoretical and empirical aspects of \ours. Methodologically, interpretability should be assessed with task-grounded and user-study metrics \citep{TonekaboniMLHC2019}, and contrasted with post-hoc explanations such as LIME and SHAP \citep[routinely used but with different guarantees]{ribeiro2016lime,lundberg2017shap}. Developing principled surrogate metrics for model selection, and procedures that reduce accuracy gaps across the simplification stages, would improve stability and reliability. Extending \ours to continuous and multi-treatment settings, as well as deriving finite-sample guarantees for CATE estimation, remain important directions. Incorporating inverse-probability weighting, doubly robust objectives, and batch sampling techniques would further broaden scope. Empirically, validation on real-world data, ideally against interventin. \add{1}{Another avenue is to adapt \ours to non-tabular modalities such as images, graphs, and genomic or sequence data by leveraging recent KAN variants for convolutional and graph architectures \citep{liu2024kan,bodner2024convolutional,bresson2024kagnn,cherednichenko2025genomics}. Designing benchmarks with ground-truth potential outcomes in these modalities is itself an open challenge, and we leave multimodal CATE benchmarking with \ours to future work.} Related with visualization, beyond PDPs, implementing related plot families, such as accumulated local effects \citep{apley2020ale} or individual conditional expectation \citep{goldstein2015peeking}, may also be informative when applied directly to the closed-form formulas. Finally, exploring robustness to adversarial perturbations, distribution shifts, and fairness constraints \citep{Rudin2018StopEB}, together with reducing the performance gaps between KAN-ification, pruning, and symbolification, will be important for reliable deployment.


